\documentclass[runningheads,a4paper]{llncs}

\usepackage[cmex10]{amsmath}
\usepackage{graphicx}
\usepackage{booktabs}
\usepackage{tabularx}
\usepackage{color}
\usepackage{subfig}
\usepackage[unicode,
					  colorlinks,
						linkcolor={red},
						citecolor={blue},
						urlcolor={blue},
            pdfauthor={Arkadiusz Jachnik, Maciej Piernik},
            pdftitle={Sequential pattern mining with uplift measure},
            pdfsubject={Sequential pattern mining},
            pdfkeywords={sequential patterns, uplift modeling}]{hyperref}

\newcolumntype{R}{>{\raggedleft\arraybackslash}X}
\newcommand{\argmax}[1]{\underset{#1}{\operatorname{argmax}}\;}
\newcommand{\matr}[1]{\mathbf{#1}}

\begin{document}

\mainmatter

\title{Sequential pattern mining with uplift measure}

\titlerunning{Sequential pattern mining with uplift measure}

\author{Arkadiusz Jachnik \and Maciej Piernik}

\authorrunning{Arkadiusz Jachnik \and Maciej Piernik}

\institute{Institute of Computing Science, Poznan University of Technology\\
ul. Piotrowo 2, 60--965 Poznan, Poland\\
\email{maciej.piernik@cs.put.poznan.pl}}

\maketitle

\begin{abstract}

\keywords{sequential patterns, uplift modeling}
\end{abstract}

\section{Introduction}
\label{sec:introduction}
Sequential data

Sequential patterns

Sequential data with status

Uplift modeling

In this paper, 

The remainder of this paper is organized as follows.

\section{Related Work}
\label{sec:related}

\section{Sequential pattern mining with uplift measure}
\label{sec:main}

\begin{figure}[!ht]
	\centering
		%\includegraphics[width=0.9\textwidth]{images/concept}
	\caption{Illustration of the proposed method}
	\label{fig:concept}
\end{figure}
\vspace{-10pt}

\section{Experiments}
\label{sec:experiments}
\subsection{Experimental Setup}

Fig.~\ref{fig:setup} depicts the experimental procedure.
\begin{figure}[h]
	\centering
		%\includegraphics[width=0.60\textwidth]{images/setup}
	\caption{Experimental procedure}
	\label{fig:setup}
\end{figure}

\subsection{Datasets}

\begin{table}[htbp]
	\scriptsize
  \centering
  \caption{Characteristics of datasets}
    \begin{tabularx}{\textwidth}{l@{}R@{}R@{}R@{}}
		\toprule
    Dataset  & \#Instances & Avg. seq. len. & \#Classes \\
		\midrule

    \bottomrule
    \end{tabularx}%
  \label{tab:datasets}%
\end{table}%

\subsection{Results}

\section{Conclusions}
\label{sec:conclusions}


\subsubsection*{Acknowledgments.} This research is partly funded by the Polish National Science Center under Grant No. DEC-2015/19/B/ST6/02637.

\bibliographystyle{splncs}
\bibliography{Bibliography}

\end{document}
